\documentclass[12pt]{article}
\usepackage[utf8]{inputenc}
\usepackage[T1]{fontenc}
\usepackage{hyperref}
\title{Amalah Language Resources for K-12 Curriculum Development}
\author{Compiled by Claude Sonnet 5.5 (claude-sonnet-5-5) for Andr\'as Kornai}
\date{}
\begin{document}
\maketitle

\section{Introduction}
Amalah (ISO 639-3 \texttt{ruy}) is also called Mala, Tumala, Amala or Rumaya. Wikipedia classes it as a Kainji language of Nigeria and gives about 6,600 speakers (a 2000 figure). Joshua Project places the Rumaya people in Lere Local Government Area, Kaduna State. Web searches found very little that could be used in a school curriculum, and this dossier lists only what turned up.

\section{Primary Educational Resources}
No primers, textbooks or school programmes for Amalah turned up in search results.

\section{Linguistic and Reference Materials}
\begin{itemize}
\item Wikipedia, ``Mala language (Nigeria)'': a short overview with classification and a speaker estimate. \url{https://en.wikipedia.org/wiki/Mala_language_(Nigeria)}
\end{itemize}
No dictionary, grammar or orthography description was found.

\section{Religious and Community Resources}
\begin{itemize}
\item Joshua Project, ``Rumaya in Nigeria'': people-group profile. It reports that a Bible translation into Amalah has been started and that audio recordings exist, but no complete New Testament. \url{https://joshuaproject.net/people_groups/14591}
\end{itemize}

\section{Digital and Audio Resources}
The Joshua Project page mentions audio recordings. No standalone audio, corpus or software resources were found.

\section{Curriculum Development Context}
The language is not known to be taught in schools, and no written educational materials were found. Any curriculum would have to start from new documentation: an orthography, a word list and recorded texts. This should be done with the community and with Bible-translation partners, who may hold unpublished materials.
\end{document}